\documentclass[11pt,a4paper]{article}
\usepackage[utf8]{inputenc}
\usepackage[french]{babel}
\usepackage{amsmath,amssymb,amsthm}
\usepackage{geometry}
\geometry{hmargin=2.5cm,vmargin=3cm}
\usepackage{tikz}
\usetikzlibrary{cd,positioning,shapes}
\usepackage{listings}
\usepackage{xcolor}

\lstset{
    language=Python,
    basicstyle=\ttfamily\small,
    keywordstyle=\color{blue}\bfseries,
    stringstyle=\color{red},
    commentstyle=\color{gray}\itshape,
    numbers=left,
    numberstyle=\tiny\color{gray},
    breaklines=true,
    frame=single,
    showstringspaces=false,
    literate={'é}{{\'e}}1 {'è}{{\`e}}1 {'à}{{\`a}}1 {'ê}{{\^e}}1 {'î}{{\^i}}1 {'ç}{{\c{c}}}1 {'ï}{{\"i}}1
}

\title{\Huge PREUVE DU THÉORÈME D'APPROXIMATION UNIVERSELLE \\ \large Cours magistral, exercices corrigés et travaux pratiques}
\author{\Large Charles EDOU NZE}
\date{\today}

\begin{document}
\maketitle
\tableofcontents
\newpage

\part{Cours Magistral}



\section{Genèse et Intuition Géométrique}

Le théorème d'approximation universelle, formalisé initialement par George Cybenko en 1989 puis étendu par Kurt Hornik en 1991, constitue le fondement mathématique justifiant la capacité des réseaux de neurones artificiels à modéliser n'importe quelle fonction continue.

L'intuition fondamentale repose sur la densité : de manière analogue au théorème d'approximation de Weierstrass qui affirme que les polynômes sont denses dans l'espace des fonctions continues sur un compact, ce théorème énonce qu'une somme pondérée de fonctions d'activation (par exemple des sigmoïdes) peut approcher uniformément toute fonction continue sur un hypercube compact de $\mathbb{R}^n$.

Géométriquement, chaque neurone d'une couche cachée agit comme un demi-plan séparateur de l'espace d'entrée. En combinant un nombre suffisant de ces demi-plans, on peut construire des "bosses" locales ou des "hyper-cylindres" qui, additionnés, permettent d'épouser n'importe quelle surface continue avec une erreur arbitrairement petite, sans nécessiter plus d'une seule couche cachée.

\begin{tikzpicture}
  \draw[->] (-0.5, 0) -- (6, 0) node[right] {$x$};
  \draw[->] (0, -0.5) -- (0, 3) node[above] {$f(x)$};

  \draw[thick, blue, domain=0:5.5, samples=100] plot (\x, {sin(\x*50) + 1.5});
  \node[blue] at (5.5, 2.5) {$f(x)$ continue};

  \draw[red, dashed, thick] (0, 1.5) -- (1, 1.5) -- (1, 2.3) -- (2, 2.3) -- (2, 0.7) -- (3, 0.7) -- (3, 1.9) -- (4, 1.9) -- (4, 1.2) -- (5, 1.2) -- (5, 0.5);
  \node[red] at (4, 2.5) {Approximation $G(x)$};
\end{tikzpicture}

\section{Définitions, Théorèmes et Exemples Concrets}

Soit $I_n = [0, 1]^n$ le cube unité de $\mathbb{R}^n$. Considérons l'espace de Banach $\mathcal{C}(I_n)$ des fonctions continues de $I_n$ dans $\mathbb{R}$, muni de la norme de la convergence uniforme :
$$ \|f\|_\infty = \sup_{x \in I_n} |f(x)| $$

\textbf{Définition (Fonction sigmoïdale)}
Une fonction $\sigma : \mathbb{R} \to \mathbb{R}$ est dite sigmoïdale si elle est mesurable et vérifie :
$$ \lim_{t \to -\infty} \sigma(t) = 0 \quad \text{et} \quad \lim_{t \to +\infty} \sigma(t) = 1 $$

\textbf{Théorème (Approximation Universelle - Cybenko 1989)}
Soit $\sigma$ une fonction d'activation continue sigmoïdale. L'ensemble $S$ des fonctions de la forme :
$$ G(x) = \sum_{j=1}^N \alpha_j \sigma(w_j^T x + b_j) $$
avec $N \in \mathbb{N}^*$, $\alpha_j, b_j \in \mathbb{R}$ et $w_j \in \mathbb{R}^n$, est dense dans $\mathcal{C}(I_n)$.
Autrement dit, pour toute $f \in \mathcal{C}(I_n)$ et tout $\epsilon > 0$, il existe $G \in S$ telle que $\|f - G\|_\infty < \epsilon$.

\textbf{Exemple 1 : Approximation d'une constante}
Si $f(x) = C$, on choisit $N=1$, $w_1 = 0$, $b_1 \to \infty$ tel que $\sigma(b_1) \approx 1$. Alors $G(x) = \alpha_1 \sigma(b_1)$. En posant $\alpha_1 = C / \sigma(b_1)$, on a $G(x) = C$.
Si $C = 5$, $\alpha_1 = 5$, $\sigma(w^T x + b) \to 1$, $G(x) \to 5$.

\textbf{Exemple 2 : Fonction porte en 1D}
Approximation de $f(x) = 1$ si $x \in [0.4, 0.6]$, $0$ sinon.
On utilise $G(x) = \sigma(k(x - 0.4)) - \sigma(k(x - 0.6))$ avec $k \gg 1$.
Pour $x = 0.5$, $k(0.1) \to +\infty \implies \sigma \to 1$, $k(-0.1) \to -\infty \implies \sigma \to 0$. $G(0.5) \approx 1 - 0 = 1$.
Pour $x = 0.2$, $k(-0.2) \to -\infty \implies \sigma \to 0$, $k(-0.4) \to -\infty \implies \sigma \to 0$. $G(0.2) \approx 0$.
Pour $x = 0.8$, $k(0.4) \to +\infty \implies 1$, $k(0.2) \to +\infty \implies 1$. $G(0.8) \approx 1 - 1 = 0$.

\textbf{Exemple 3 : Évaluation numérique de la porte}
Avec $\sigma(t) = \frac{1}{1+e^{-t}}$, pour $k=100$ et $x=0.5$:
$G(0.5) = \sigma(10) - \sigma(-10) = \frac{1}{1+e^{-10}} - \frac{1}{1+e^{10}} \approx 0.99995 - 0.00005 = 0.9999$.

\textbf{Exemple 4 : La bosse 2D}
Soit $x = (x_1, x_2) \in [0,1]^2$. Pour approcher un pic autour de $(0.5, 0.5)$, on somme plusieurs "hyper-cylindres" orientés le long de différents vecteurs $w$.
Si on veut localiser l'activation, on prend $\sigma(w_1 x_1 - b_1) - \sigma(w_1 x_1 - b_2)$ et similairement pour $x_2$.

\textbf{Exemple 5 : Fonction d'activation ReLU}
Le théorème se généralise pour $\sigma(t) = \max(0, t)$. Une fonction "chapeau" s'écrit: $T(x) = \sigma(x+1) - 2\sigma(x) + \sigma(x-1)$.
Pour $x=0$, $T(0) = \sigma(1) - 2\sigma(0) + \sigma(-1) = 1 - 0 + 0 = 1$.
Pour $x=1$, $T(1) = \sigma(2) - 2\sigma(1) + \sigma(0) = 2 - 2 + 0 = 0$.
Pour $x=-1$, $T(-1) = \sigma(0) - 2\sigma(-1) + \sigma(-2) = 0 - 0 + 0 = 0$.

\textbf{Exemple 6 : Limites pathologiques}
Si $f(x) = \sin(1/x)$ sur $(0, 1]$ prolongée par 0 en 0. $f$ n'est pas continue en 0. Le théorème ne s'applique pas sur le compact $[0,1]$ au sens de la norme uniforme.

\begin{tikzpicture}
  \draw[->] (-2, 0) -- (2, 0) node[right] {$x$};
  \draw[->] (0, -0.5) -- (0, 2) node[above] {$\sigma(x)$};

  \draw[thick, blue, domain=-2:2, samples=100] plot (\x, {max(0, \x)});
  \node[blue] at (1.5, 1.8) {ReLU};

  \draw[thick, red, domain=-2:2, samples=100] plot (\x, {1 / (1 + exp(-4*\x))});
  \node[red] at (-1, 0.5) {Sigmoïde};
\end{tikzpicture}


\section{Démonstrations}

La preuve originale exploite la dualité topologique et le théorème de Riesz-Markov, démontrant la densité par l'absurde via le théorème de Hahn-Banach.

\textbf{Étape 1 : Hypothèse par l'absurde}
Soit $S$ le sous-espace vectoriel engendré par les fonctions $x \mapsto \sigma(w^T x + b)$.
Supposons que $S$ ne soit pas dense dans $\mathcal{C}(I_n)$.
L'adhérence $\bar{S}$ est donc un sous-espace fermé strict de $\mathcal{C}(I_n)$.

\textbf{Étape 2 : Intervention du Théorème de Hahn-Banach}
D'après un corollaire du théorème de Hahn-Banach (forme analytique), il existe une forme linéaire continue non nulle $L$ sur $\mathcal{C}(I_n)$ telle que $L(g) = 0$ pour toute fonction $g \in \bar{S}$.

\textbf{Étape 3 : Théorème de Représentation de Riesz}
Le dual de $\mathcal{C}(I_n)$ est isomorphe à l'espace des mesures de Borel régulières signées finies sur $I_n$.
Il existe donc une mesure $\mu \neq 0$ sur $I_n$ telle que pour toute $f \in \mathcal{C}(I_n)$, $L(f) = \int_{I_n} f(x) d\mu(x)$.
Puisque $L$ s'annule sur $S$, on a :
$$ \int_{I_n} \sigma(w^T x + b) d\mu(x) = 0 \quad \forall w \in \mathbb{R}^n, \forall b \in \mathbb{R} $$

\textbf{Étape 4 : Propriété discriminatoire}
Une fonction $\sigma$ est dite discriminatoire si la condition ci-dessus implique $\mu = 0$.
Cybenko démontre que toute fonction sigmoïdale continue est discriminatoire. En fixant $w \neq 0$, et en considérant $h(x) = w^T x$, on projette la mesure $\mu$ sur la droite engendrée par $w$.
Par des arguments d'analyse de Fourier (ou par transformation de Laplace pour des mesures à support compact), l'annulation de l'intégrale pour toutes translations $b$ et dilatations des sigmoïdes implique que la mesure projetée est nulle pour toute direction $w$.

\textbf{Étape 5 : Conclusion}
Puisque les projections monodimensionnelles de la mesure $\mu$ sont toutes nulles, la transformée de Fourier de $\mu$ s'annule partout. Cela entraîne que $\mu$ est la mesure nulle.
Or, on avait supposé $L \neq 0$, donc $\mu \neq 0$, ce qui est une contradiction.
Par conséquent, $S$ est nécessairement dense dans $\mathcal{C}(I_n)$.

\section{Applications en Physique, Logique et IA}

Le théorème d'approximation universelle justifie théoriquement la capacité de l'apprentissage profond à modéliser des systèmes physiques complexes non linéaires, tels que la dynamique des fluides (via les Physics-Informed Neural Networks - PINNs) ou la résolution d'équations aux dérivées partielles à haute dimension (comme l'équation de Schrödinger).

En logique, ce théorème établit que les réseaux de neurones constituent un système formel complet pour l'approximation de fonctions continues réelles, surpassant les limites expressives du perceptron simple, qui ne pouvait pas représenter la fonction logique XOR non linéairement séparable.


\newpage
\part{Exercices d'Application et de Concours}
\subsection*{Exercice 1 : Approximation d'une constante $\bigstar$}
On considère le cube $I_1 = [0,1]$ et la fonction d'activation sigmoïde $\sigma(t) = \frac{1}{1+e^{-t}}$.
Soit $f(x) = 3$ pour tout $x \in I_1$.
Construire explicitement une fonction $G(x) = \alpha \sigma(wx + b)$ telle que $\|f - G\|_\infty < 0.01$.

\textbf{Correction détaillée}
On cherche $G(x) = \alpha \sigma(wx + b)$.
Posons $w = 0$, la fonction devient constante par rapport à $x$ : $G(x) = \alpha \sigma(b)$.
On veut approcher la valeur 3. Si on choisit $b$ très grand, par exemple $b = 10$, on a $\sigma(10) = \frac{1}{1+e^{-10}} \approx 0.9999546$.
Pour que $G(x) = 3$, il faut choisir $\alpha = \frac{3}{\sigma(10)} = 3(1+e^{-10}) \approx 3.000136$.
Ainsi, avec $w=0, b=10, \alpha = 3(1+e^{-10})$, on a $G(x) = 3$ exactement pour tout $x \in [0,1]$.
L'erreur est rigoureusement nulle, ce qui satisfait $\|f - G\|_\infty < 0.01$.


\subsection*{Exercice 2 : La porte logique NOT $\bigstar\star\star\star\star$}
Approcher la fonction $f(x) = 1 - x$ sur $[0,1]$ avec un réseau à une couche cachée utilisant l'activation linéaire par morceaux $\sigma(x) = \max(0, \min(1, x))$ (activation Hard Sigmoid).

\textbf{Correction détaillée}
La fonction cible est affine : $f(x) = -x + 1$.
La fonction d'activation est $\sigma(t) = t$ pour $t \in [0,1]$.
Si on choisit l'argument $w x + b$ tel qu'il reste dans $[0,1]$ pour tout $x \in [0,1]$.
Posons $w = 1, b = 0$, on a $\sigma(x) = x$ pour $x \in [0,1]$.
On cherche $G(x) = \alpha \sigma(wx+b) + c$ où on autorise un biais de sortie.
Le théorème d'approximation dit qu'on peut utiliser une combinaison linéaire : $G(x) = \alpha_1 \sigma(w_1 x + b_1) + \alpha_2 \sigma(w_2 x + b_2)$.
En choisissant $w_1 = 1, b_1 = 0$, on a $\sigma(x)$.
En choisissant $w_2 = 0, b_2 = 1$, on a $\sigma(1) = 1$.
Ainsi, $G(x) = -\sigma(x) + \sigma(1) = -x + 1$.
La fonction est modélisée exactement avec deux neurones.


\subsection*{Exercice 3 : La porte logique XOR (version continue) $\bigstar\bigstar\star\star\star$}
Considérons $I_2 = [0,1]^2$. On veut approcher la fonction surface $f(x_1, x_2) = x_1 + x_2 - 2x_1 x_2$.
Montrer comment utiliser des ReLUs $\sigma(t) = \max(0, t)$ pour approcher le terme croisé $x_1 x_2$.

\textbf{Correction détaillée}
On sait que $2x_1 x_2 = (x_1 + x_2)^2 - x_1^2 - x_2^2$.
Pour approcher la fonction carré $t \mapsto t^2$ sur $[0, 2]$, on peut utiliser une somme de ReLUs.
Puisque $t^2$ est convexe, on peut l'approcher par une ligne brisée convexe, qui est une combinaison linéaire de fonctions $\max(0, t - t_i)$.
Soit $N$ le nombre de segments. On subdivise $[0, 2]$ en $t_i = i \frac{2}{N}$.
La fonction approchant $t^2$ est de la forme $H_N(t) = \sum_{i=1}^{N-1} c_i \max(0, t - t_i) + c_0 t + b_0$.
Ainsi, la multiplication $x_1 x_2 = \frac{1}{2} (H_N(x_1+x_2) - H_N(x_1) - H_N(x_2))$ peut être approchée arbitrairement près.
Puisque $H_N$ est une somme de ReLUs, $x_1 x_2$ s'exprime comme une combinaison linéaire de ReLUs de la forme $\sigma(w^T x + b)$.
La densité s'ensuit.


\subsection*{Exercice 4 : Propriété Discriminatoire $\bigstar\bigstar\bigstar\star\star$}
Soit $\mu$ une mesure signée finie sur $\mathbb{R}$. Supposons que $\int_{\mathbb{R}} \exp(i w x) d\mu(x) = 0$ pour tout $w \in \mathbb{R}$. En déduire que $\mu = 0$.

\textbf{Correction détaillée}
L'hypothèse indique que la transformée de Fourier de la mesure $\mu$, notée $\mathcal{F}(\mu)(w) = \int_{\mathbb{R}} e^{-iwx} d\mu(x)$, est identiquement nulle pour tout $w \in \mathbb{R}$.
L'espace des mesures régulières signées finies sur $\mathbb{R}$ correspond au dual de l'espace des fonctions continues s'annulant à l'infini $C_0(\mathbb{R})$.
L'application qui à une mesure associe sa transformée de Fourier est injective. (Théorème d'inversion de Fourier étendu aux mesures, ou densité de l'algèbre engendrée par les caractères exponentiels via Stone-Weierstrass complexe sur la compactification).
Puisque $\mathcal{F}(\mu) = 0$, on conclut immédiatement que $\mu = 0$.
Cette propriété est cruciale dans la preuve de Cybenko lorsque l'on exprime la sigmoïde comme une superposition de fonctions exponentielles complexes (via la transformée de Fourier).


\subsection*{Exercice 5 : Densité des fonctions en escalier $\bigstar\bigstar\bigstar\star\star$}
Montrer que l'ensemble des fonctions en escalier sur $I_1 = [0,1]$ est dense dans $L^1([0,1])$.

\textbf{Correction détaillée}
Par définition de l'intégrale de Lebesgue, les fonctions étagées (combinaisons linéaires d'indicatrices d'ensembles mesurables) sont denses dans $L^1$.
Un sous-ensemble mesurable $A$ peut être approché en mesure par des unions finies d'intervalles ouverts.
L'indicatrice d'un intervalle est une fonction en escalier.
Donc, pour toute fonction intégrable $f \in L^1([0,1])$ et $\epsilon > 0$, il existe une fonction étagée $\phi = \sum_{i=1}^n c_i \mathbb{I}_{A_i}$ telle que $\|f - \phi\|_1 < \epsilon/2$.
Ensuite, pour chaque $A_i$, il existe une union finie d'intervalles $U_i$ telle que $m(A_i \Delta U_i) < \epsilon / (2n|c_i|)$.
La fonction $\psi = \sum_{i=1}^n c_i \mathbb{I}_{U_i}$ est en escalier.
Par l'inégalité triangulaire, $\|f - \psi\|_1 \le \|f - \phi\|_1 + \|\phi - \psi\|_1 \le \epsilon/2 + \sum |c_i| m(A_i \Delta U_i) < \epsilon$.


\subsection*{Exercice 6 : Non-linéarité stricte $\bigstar\bigstar\bigstar\bigstar\star$}
Si la fonction d'activation est $\sigma(t) = at + b$, le théorème d'approximation universelle est-il valable ?

\textbf{Correction détaillée}
Non. L'ensemble $S$ des réseaux s'écrit $G(x) = \sum_{j=1}^N \alpha_j (a(w_j^T x) + b_j + b) = (\sum \alpha_j a w_j)^T x + \sum \alpha_j(b_j+b)$.
C'est-à-dire $G(x) = W^T x + B$ où $W \in \mathbb{R}^n$ et $B \in \mathbb{R}$.
$S$ est exactement l'espace des fonctions affines sur $\mathbb{R}^n$.
L'espace des fonctions affines est un sous-espace vectoriel de dimension finie $n+1$.
Dans l'espace $\mathcal{C}(I_n)$ qui est de dimension infinie, un sous-espace de dimension finie est toujours fermé (et complet), et ne peut donc pas être dense.
Il est impossible d'approcher uniformément une fonction non-affine, par exemple $f(x) = x^2$ sur $[0,1]$.
Si $\sup_{x \in [0,1]} |x^2 - (Wx+B)| < \epsilon$, alors les dérivées secondes devraient être proches, ce qui n'est pas le cas.


\subsection*{Exercice 7 : Approximation d'une gaussienne $\bigstar\bigstar\bigstar\bigstar\star$}
Approcher $f(x) = \exp(-x^2)$ sur $[-1, 1]$ par une combinaison de ReLU $\max(0, x)$.

\textbf{Correction détaillée}
$f$ est dérivable et convexe par morceaux (inflexion en $\pm 1/\sqrt{2}$).
On utilise une subdivision de $[-1, 1]$ avec pas $h = 2/N$, $x_i = -1 + ih$.
L'interpolée linéaire par morceaux est $G(x)$. La dérivée seconde de $f$ est bornée par 2, donc l'erreur d'interpolation spline linéaire est encadrée par $M_2 h^2 / 8 \le 2 / (8 N^2) = 1/(4N^2)$.
Pour avoir une erreur $\epsilon$, il suffit que $1/(4N^2) < \epsilon$, soit $N > \frac{1}{2\sqrt{\epsilon}}$.
La fonction spline $G(x)$ s'écrit formelently $G(x) = f(-1) + c_{-1}(x+1) + \sum_{i=1}^{N-1} a_i \max(0, x-x_i)$.
Où les coefficients $a_i$ représentent les sauts de dérivée aux noeuds $x_i$ : $a_i = f'(x_i+) - f'(x_i-)$.
Ainsi, l'approximation est bien une somme pondérée de ReLUs.


\subsection*{Exercice 8 : Théorème de Stone-Weierstrass $\bigstar\bigstar\bigstar\bigstar\bigstar$}
En utilisant le théorème de Stone-Weierstrass, démontrer l'approximation universelle avec $\sigma(x) = x^2$ pour un domaine compact arbitraire $K \subset \mathbb{R}^n$.

\textbf{Correction détaillée}
Considérons l'ensemble $A = \text{Vect}\{ (w^T x + b)^2 \mid w \in \mathbb{R}^n, b \in \mathbb{R} \}$.
Développons le carré : $(w^T x + b)^2 = (w^T x)^2 + 2b w^T x + b^2 = \sum_{i,j} w_i w_j x_i x_j + 2b \sum w_i x_i + b^2$.
On voit que $A$ contient :
- Les constantes (en posant $w=0, b \neq 0$).
- Les fonctions affines $x_i$ (en ajustant $b$).
- Les termes quadratiques et croisés $x_i^2$ et $x_i x_j$ (par polarisation $(x_i+x_j)^2 - x_i^2 - x_j^2$).
Toutefois, l'algèbre engendrée par $A$ engendre tous les polynômes de $\mathbb{R}^n$, qui est dense dans $\mathcal{C}(K)$ par Stone-Weierstrass.
Cependant, l'ensemble $A$ lui-même ne contient QUE les polynômes de degré au plus 2.
Le théorème d'approximation universelle stricto sensu n'est pas vérifié pour les réseaux à une SEULE couche cachée avec $\sigma(x)=x^2$, car ils ne peuvent générer des degrés $\ge 3$.
L'affirmation est fausse pour une seule couche. Elle nécessite plusieurs couches (fonctions composées) pour que $\sigma(x)=x^2$ puisse engendrer par compositions successives des polynômes de degrés arbitrairement élevés, rendant le réseau dense.


\subsection*{Exercice 9 : Hahn-Banach et Dualité $\bigstar\bigstar\bigstar\bigstar\bigstar$}
Justifier rigoureusement que si $\bar{S} \neq E$ (où $E$ est un espace de Banach réel et $S$ un sous-espace vectoriel), il existe une forme linéaire continue $L \in E^*$ telle que $L_{|S} = 0$ et $L \neq 0$.

\textbf{Correction détaillée}
Soit $x_0 \in E \setminus \bar{S}$.
Puisque $\bar{S}$ est un sous-espace vectoriel fermé, la distance $d = \inf_{y \in \bar{S}} \|x_0 - y\|$ est strictement positive ($d > 0$).
Considérons le sous-espace $F = \bar{S} \oplus \mathbb{R}x_0$.
Définissons une forme linéaire $f_0 : F \to \mathbb{R}$ par $f_0(y + \lambda x_0) = \lambda d$ pour tout $y \in \bar{S}$ et $\lambda \in \mathbb{R}$.
Vérifions sa continuité sur $F$ :
Pour $\lambda \neq 0$, $\|y + \lambda x_0\| = |\lambda| \| \frac{y}{\lambda} + x_0 \| \ge |\lambda| d$.
Donc $|f_0(y + \lambda x_0)| = |\lambda| d \le \|y + \lambda x_0\|$.
Ainsi, $f_0$ est continue et sa norme sur $F$ est $\le 1$ (en fait exactement 1).
Par le théorème de Hahn-Banach analytique, il existe un prolongement linéaire $L : E \to \mathbb{R}$ tel que $L_{|F} = f_0$ et $\|L\|_{E^*} = \|f_0\|_{F^*} = 1$.
Puisque pour $y \in S$, $L(y) = f_0(y) = 0$, $L$ s'annule sur $S$.
Comme $L(x_0) = f_0(x_0) = d > 0$, $L$ n'est pas la forme nulle.


\subsection*{Exercice 10 : Mesure de Radon $\bigstar\bigstar\bigstar\bigstar\bigstar$}
Soit $I_n$ compact. Si $\int_{I_n} f(x) d\mu(x) = 0$ pour toute $f \in \mathcal{C}(I_n)$, démontrer que la mesure de Radon signée $\mu$ est nulle, en invoquant le théorème de représentation de Riesz.

\textbf{Correction détaillée}
Le théorème de représentation de Riesz-Markov-Kakutani énonce que pour tout espace localement compact séparé $X$ (ici $I_n$ compact, donc localement compact), il existe un isomorphisme isométrique entre l'espace dual $C_0(X)^*$ (ici $\mathcal{C}(I_n)^*$ car $I_n$ compact) et l'espace des mesures de Radon régulières signées finies sur $X$, noté $\mathcal{M}(X)$, muni de la norme en variation totale.
Soit l'opérateur $T : \mathcal{M}(I_n) \to \mathcal{C}(I_n)^*$ défini par $T(\mu)(f) = \int_{I_n} f d\mu$.
Le fait que $T$ soit une isométrie implique que $\|T(\mu)\|_{C^*} = \|\mu\|_{TV}$.
L'hypothèse indique que pour toute $f \in \mathcal{C}(I_n)$, $T(\mu)(f) = 0$.
Donc la forme linéaire $T(\mu)$ est la forme identiquement nulle dans $\mathcal{C}(I_n)^*$.
Par l'isométrie, $\|\mu\|_{TV} = \|0\|_{C^*} = 0$.
Une mesure dont la variation totale est nulle est identiquement nulle sur tous les boréliens. Donc $\mu = 0$.




\newpage
\part{Travaux Pratiques et Simulations Algorithmiques}
\subsection*{TP 1 : Approximation d'une fonction sinusoïdale 1D}


\begin{lstlisting}
# Pure Python implementation
import math
import random

def sigmoid(x):
    return 1.0 / (1.0 + math.exp(-x))

def relu(x):
    return max(0.0, x)

def forward(x, weights, biases, v_weights, v_bias):
    # Hidden layer
    hidden = []
    for i in range(len(weights)):
        hidden.append(sigmoid(weights[i] * x + biases[i]))

    # Output layer
    out = v_bias
    for i in range(len(v_weights)):
        out += v_weights[i] * hidden[i]
    return out

def mse(y_true, y_pred):
    return (y_true - y_pred)**2

# Dataset: f(x) = sin(x) on [0, 2*pi]
x_data = [i * 2.0 * math.pi / 100.0 for i in range(101)]
y_data = [math.sin(x) for x in x_data]

# Initialize parameters for N=10 neurons
N = 10
weights = [random.uniform(-1, 1) for _ in range(N)]
biases = [random.uniform(-1, 1) for _ in range(N)]
v_weights = [random.uniform(-1, 1) for _ in range(N)]
v_bias = 0.0

# Simple SGD
lr = 0.01
epochs = 1000
for epoch in range(epochs):
    total_loss = 0
    for x, y in zip(x_data, y_data):
        # Forward
        hidden = [sigmoid(weights[i] * x + biases[i]) for i in range(N)]
        y_pred = v_bias + sum(v_weights[i] * hidden[i] for i in range(N))

        # Loss
        loss = (y - y_pred)**2
        total_loss += loss

        # Backward (gradient calculation)
        grad_y_pred = -2 * (y - y_pred)

        v_bias -= lr * grad_y_pred
        for i in range(N):
            grad_vw = grad_y_pred * hidden[i]
            grad_h = grad_y_pred * v_weights[i]

            # sigmoid derivative: s(x)*(1-s(x))
            grad_z = grad_h * hidden[i] * (1 - hidden[i])

            grad_w = grad_z * x
            grad_b = grad_z * 1

            v_weights[i] -= lr * grad_vw
            weights[i] -= lr * grad_w
            biases[i] -= lr * grad_b

    # Mathematical validation assert
    assert total_loss >= 0, "Loss cannot be negative"

print("Final Loss:", total_loss / len(x_data))
\end{lstlisting}



\subsection*{TP 2 : Construction de la fonction porte logique (Step Function)}


\begin{lstlisting}
# Pure Python implementation
import math

def sigmoid(x):
    return 1.0 / (1.0 + math.exp(-x))

# We want to approximate f(x) = 1 if 0.4 <= x <= 0.6 else 0
# using G(x) = sigmoid(k*(x-0.4)) - sigmoid(k*(x-0.6))

def approximate_gate(x, k):
    return sigmoid(k * (x - 0.4)) - sigmoid(k * (x - 0.6))

test_points = [0.0, 0.2, 0.4, 0.5, 0.6, 0.8, 1.0]

for k in [10, 100, 1000]:
    print("Results for k =", k)
    for x in test_points:
        val = approximate_gate(x, k)
        print("x =", x, "=> G(x) =", val)
    print("---")

# Mathematical validation assert
val_center = approximate_gate(0.5, 1000)
val_edge = approximate_gate(0.2, 1000)
assert val_center > 0.99, "Should be close to 1"
assert val_edge < 0.01, "Should be close to 0"
\end{lstlisting}



\subsection*{TP 3 : Approximation avec activation ReLU}


\begin{lstlisting}
# Pure Python implementation

def relu(x):
    return x if x > 0 else 0.0

# Approximate f(x) = x^2 on [0, 2] using 4 segments
# Interpolation nodes: 0, 0.5, 1.0, 1.5, 2.0
# True values: 0, 0.25, 1.0, 2.25, 4.0
# Slopes of segments: 0.5, 1.5, 2.5, 3.5

def approx_square(x):
    # Form: f(0) + c0*x + sum_i a_i relu(x - xi)
    # a_i = change in slope at xi

    slope0 = 0.5
    slope1 = 1.5  # change = 1.0 at x=0.5
    slope2 = 2.5  # change = 1.0 at x=1.0
    slope3 = 3.5  # change = 1.0 at x=1.5

    val = slope0 * x
    val += 1.0 * relu(x - 0.5)
    val += 1.0 * relu(x - 1.0)
    val += 1.0 * relu(x - 1.5)

    return val

test_points = [0.0, 0.5, 1.0, 1.5, 2.0]
for x in test_points:
    true_val = x**2
    approx = approx_square(x)
    print("x:", x, "True:", true_val, "Approx:", approx)

    # Mathematical validation assert
    assert abs(true_val - approx) < 1e-9, "Exact interpolation at nodes expected"
\end{lstlisting}



\subsection*{TP 4 : Transformée de Fourier empirique d'une Sigmoïde}


\begin{lstlisting}
# Pure Python implementation
import math

def sigmoid(x):
    return 1.0 / (1.0 + math.exp(-x))

# We want to approximate the integral of sigmoid(x) * exp(-i*w*x) dx
# using a Riemann sum on [-L, L]

def empirical_fourier_transform(w, L, N):
    dx = 2.0 * L / N
    integral_real = 0.0
    integral_imag = 0.0

    for i in range(N):
        x = -L + i * dx
        val = sigmoid(x)
        integral_real += val * math.cos(w * x) * dx
        integral_imag -= val * math.sin(w * x) * dx

    return integral_real, integral_imag

w_vals = [0.0, 1.0, 5.0]
L = 10.0
N = 1000

for w in w_vals:
    r, i = empirical_fourier_transform(w, L, N)
    print("w:", w, "Real:", r, "Imag:", i)

# Mathematical validation assert
# for w=0, integral of sigmoid on [-L, L] should be L
r0, _ = empirical_fourier_transform(0.0, L, N)
assert abs(r0 - L) < 1.0, "Integral of sigmoid near 0 over symmetric domain is approx domain half width"
\end{lstlisting}



\subsection*{TP 5 : Limite de la profondeur vs largeur (Démonstrateur)}


\begin{lstlisting}
# Pure Python implementation
import math
import random

def relu(x):
    return x if x > 0 else 0.0

# Compare modeling parity function (XOR like) with 1 hidden layer vs 2 hidden layers
# Parity of 3 bits: 1 if sum of bits is odd, 0 if even

def generate_parity_data():
    data = []
    for i in range(8):
        b1 = i & 1
        b2 = (i >> 1) & 1
        b3 = (i >> 2) & 1
        y = (b1 + b2 + b3) % 2
        data.append(([float(b1), float(b2), float(b3)], float(y)))
    return data

def forward_1_layer(x, w, b, vw, vb):
    h = [relu(sum(w[i][j]*x[j] for j in range(3)) + b[i]) for i in range(len(b))]
    return sum(vw[i]*h[i] for i in range(len(b))) + vb

data = generate_parity_data()

# 1 layer with N neurons
N = 4 # Minimal required for 3D XOR is 4
w = [[random.uniform(-1,1) for _ in range(3)] for _ in range(N)]
b = [random.uniform(-1,1) for _ in range(N)]
vw = [random.uniform(-1,1) for _ in range(N)]
vb = 0.0

# Simple numerical perturbation for gradient check / learning demo
lr = 0.01
# Training loop omitted for brevity, but mathematically it's proven N=4 is sufficient.

# Assert dataset size
assert len(data) == 8
\end{lstlisting}





\end{document}
